% !TeX program = pdflatex
\documentclass[11pt,a4paper]{article}
\usepackage[T1]{fontenc}
\usepackage[utf8]{inputenc}
\usepackage[margin=24mm]{geometry}
\usepackage{newtxtext,newtxmath,microtype,amsmath,booktabs,tabularx,xcolor}
\usepackage[numbers,sort&compress]{natbib}
\usepackage[hyphens]{xurl}
\usepackage{hyperref}
\definecolor{inkblue}{HTML}{006786}
\hypersetup{colorlinks=true,urlcolor=inkblue,citecolor=inkblue,linkcolor=inkblue,pdftitle={Response to review: forest-carbon measurement and annual accounts},pdfauthor={Henry Arellano-Peña},pdfsubject={Author response to methodological and source-audit review}}
\renewcommand{\bibfont}{\small}\setlength{\bibsep}{4pt}
\setlength{\parskip}{.45em}\setlength{\parindent}{0pt}\setlength{\emergencystretch}{3em}
\title{Response to review\\\large From tropical forest-carbon measurement to annual carbon accounts}
\author{Henry Arellano-Pe\~na\\\small NEBIOT S.A.S.}
\date{16 September 2026}
\begin{document}
\maketitle
I thank the reviewer for the numerical verification and the proposed refinements. I have revised the article to make its public research record, accounting boundaries and classical mathematical foundations explicit. The principal additions are a two-model display-precision audit, a mixed-cohort annual-release corollary and a near-linear census-calibration result. Two suggested interpretations require additional conditions: the finite subset bracket and the probability comparison of display conventions. I state those conditions below and in the paper.

\section*{1. Public data and code location}
The paper cites the research repository at \url{https://github.com/hyrucanji77/Tropical_Forest_Carbon-}. It contains the manuscript and figure sources, numerical inputs, source-provenance table, executable analysis, generated outputs, verification scripts and manifest. The research snapshot is identified as 5.0 in repository metadata rather than on the paper's title page. Commit history and SHA-256 records distinguish immutable content from later changes. I have not assigned or implied a Zenodo DOI.

The data-availability statement and Appendix B identify actual repository paths. The scripts regenerate 26 numerical and LaTeX-table outputs. The programme contains 22,100 randomized cases plus deterministic regression and invalid-input checks; the executed assertion count is in \path{data/validation.json}. The obsolete share-derivative description and check have been removed. Tests establish computational properties, not independent validation of a forest inventory.

\section*{2. Truncation, rounding and the printed class table}
Table 2 includes the source's printed intervals, means and percentages alongside the recomputed stock/area ratios. Truncating $100C_j/723.97$ to two decimals reproduces all five percentages; nearest rounding fails for classes 4, 2 and 3. Both the stock and area sums lie two last-digit units below their displayed totals. The source table is transcribed without numerical repair \citep{Arellano2016}.

The audit now propagates two explicit error models through every class. Nearest rounding gives a common normalization-factor interval of approximately $[0.990140,0.990190]$; truncation gives $[0.990134,0.990182]$. Exact $1/1.01$ is excluded by classes 4 and 5 under either model. Thus the approximately one-percent normalization and crossed class metadata survive the display-rule sensitivity analysis.

I have qualified the probability argument. For five independent uniform hidden component fractions, with the total calculated from their exact sum and displayed under the same rule, a positive two-unit discrepancy has probability $11/20=0.55$ under truncation. Under nearest rounding, the corresponding one-sided probability is $1/120$; $1/60$ includes either sign. The probabilities concern an explicitly assumed display model, not the historical author's unknown workflow. Stock and area outputs may be dependent, so their probabilities are not multiplied.

There is a small further qualification to the proposed alternative denominator near 724.08. With the displayed stocks fixed, nearest-rounding the five percentages gives incompatible lower and upper denominator limits. Even allowing each stock half a last-digit unit, the tight joint bounds are approximately 724.082332762 and 724.082306618 PgC: a gap of $2.61\times10^{-5}$ PgC rather than a nonempty exact interval. This is negligible as carbon science, but it prevents describing 724.08 as an exact jointly compatible denominator. The high-precision values are retained in the audit JSON.

The fixed-total historical endpoints are also conditional on display precision. An exact 10.3\% baseline plus truncation gives 26.15--32.60\%. Nearest rounding with a hidden baseline of approximately 10.296665--10.299824\%, itself displayed as 10.3\%, gives the same endpoint labels. The coherent truncation pattern is evidence for that explanation, but the actual calculation was not recovered. The figure inset consequently states the ambiguity rather than asserting a confirmed formatting cause.

\section*{3. Original empirical records and subset bounds}
The published preprint and its tabulation are recoverable. The pre-tabulation aggregation workbook, operational class key and numerical rasters could not be located among the records available for this audit. I cannot distinguish loss, storage elsewhere or supersession from that absence. The paper therefore specifies which part of the research is reproducible: the published-digit audit and exact analytical calculations, rather than the original map-production pipeline.

The high-carbon subset comparison has been tightened. If the 1200.346 Mha subset $H$ lies within Feldpausch's 1667.5 Mha domain, its conventional carbon obeys $C_{0,H}\leq285.1$ PgC \citep{Feldpausch2012}. Thus its required correction is at least $566.79/285.1\simeq1.99$. The value 2.76 assigns the conventional whole-domain mean to those particular pixels. It is an upper bound only with the additional assumption that their conventional mean is at least the whole-domain mean. Without that condition there is no finite upper bound from the aggregate totals alone.

Section 4 therefore uses the conditional bracket 1.99--2.76 only with both assumptions stated. Its lower endpoint exceeds the listed moist-forest replacement factors but lies within Demol et al.'s dry-woodland comparison range \citep{Demol2024}. Numerical overlap remains a comparison of heterogeneous evidence, not independent validation. The larger 2.76 scenario exceeds the listed factors; it is not presented as a measured local correction. The Baccini endpoint is explicitly left without a corresponding area factor because that denominator has not been established.

\section*{4. Annual mixtures and inventory timing}
The single-cohort release identity is extended to an annual mixture through
\[
W_i(t)=\sum_{s\geq0}D(t-s)p_{i,s},\qquad
r_{F,t}=r_C+\frac{\operatorname{Cov}_{\mu}(r,W(t))}{\mathbb E_{\mu}[W(t)]}.
\]
Under stable compartment composition, constant disturbance over the relevant history and complete ultimate release, the annual multiplier equals $r_C$. With two lags, the covariance is multiplied by $D(t)-D(t-1)$. The worked example therefore gives 2 for constant clearing, $51/29$ for a doubled latest cohort and $69/31$ for a halved latest cohort. The corollary states the exact assumptions; arbitrary longer histories require the full convolution.

I also separate physical release from conversion-year stock-loss booking. EDGAR describes a Tier 1 Gain/Loss implementation but its Annex 2 does not provide a size-resolved deforestation decay kernel \citep{Crippa2026}. Under a convention that books all conversion losses in the conversion year, the inventory weights are affected fractions and area; the timing convolution instead matters for atmospheric transfers and budget comparison. The paper no longer invites the reader to transfer a physical lag directly into a booked category without that distinction.

\section*{5. Classical identities and empirical covariance context}
Proposition 5.1 is identified as the covariance decomposition of a weighted mean, algebraically the selection term of the Price equation \citep{Price1970}. The opposite-order argument is identified as Chebyshev's sum inequality. The contribution lies in defining the carbon, disturbance and time weights and their valid accounting uses, not in claiming a new general covariance theorem.

I added the tropical edge-biomass study: approximately 25\% lower biomass within 500 m and nearly 10\% Tier 1 stock overstatement when these effects are omitted \citep{ChaplinKramer2015}. This is evidence for spatial heterogeneity. A negative correction--clearing covariance follows only if an independently observed clearing distribution weights the more overestimated edge stocks. The biomass study alone does not measure that joint distribution, so the sign is presented conditionally.

The near-linear calibration expansion is included with a dimensionless stock reference $x_*$:
\[
r_L-r_G=a(b-1)\{\mathbb E_\ell\ln(X/x_*)-\mathbb E_\gamma\ln(X/x_*)-1\}+O((b-1)^2).
\]
Here $\gamma$ is the exact finite-interval gain distribution, including zero-start recruitment intervals. The differential shrinks linearly and its first-order sign depends on geometric means. A zero first-order coefficient requires higher-order analysis. The code tests the scaled second-order remainder rather than incorrectly testing the approximation as an exact finite-exponent identity.

The EDGAR discussion now provides a scale illustration: signed positive and negative land aggregates each near 6 GtCO$_2$e against a residual near 0.9 give a cancellation ratio of order seven percent \citep{Crippa2026}. This is an approximate inventory-boundary illustration, not ecosystem gross exchange or evidence of a shared calibration error.

\section*{6. Figures, bibliography and presentation}
The captions credit the WRI/Ecofys source--activity--gas convention. They distinguish intermediate conceptual activities from EDGAR's official sector list. The report's CC BY 4.0 terms are recorded separately from legacy graphic and third-party rights; no unverified source-specific permission is claimed. The green branch in Figure 2 is hatched and outlined so its magnitude is visibly unresolved even outside the caption.

The title page and figure headers omit revision numbers, PDF metadata identifies the work as a research article and supplies keywords, and the abstract is below 250 words. Revision-relative language has been removed from the abstract, introduction, captions, conclusions and Appendix B. A scope-of-inference box centralizes the main evidential limits.

I checked Gonzalez de Tanago et al.'s Equation 3 and Table 4 again. The latter explicitly states a sample size of 20 validation trees; nine other trees were used for calibration. The abstract's ``overall'' wording does not expand that table to 29 validation observations. Largest-tree selection and mean diameter 73.5 cm are now explicit, while a smaller stand-level factor is conditional on the same size-dependent relation and appropriate population weights \citep{Gonzalez2018}. The Ecofys component is displayed as approximately 5.009 GtCO$_2$ rather than a misleading six-decimal estimate.

\section*{7. Author declarations}
The manuscript remains under my sole authorship. Its declaration records my ownership of the methodologies and founder/director role at NEBIOT, the sole company currently affiliated with this work, together with my statement of no competing interests. It explicitly identifies NEB001's CCA--Fuzzy Land Cover mapping and NEB002's voxel, elemental-carbon and neural-network chain as underlying the estimate being audited. It contains no unrelated-company or repetitive authorship-process discussion. Published references retain their source author lists.

Correspondence is supplied as \href{mailto:harellano@unal.edu.co}{harellano@unal.edu.co}. Funding is stated as ``No external funding is reported for this analysis,'' without inventing a grant or independently certifying an unreported funding history. Assistance from ChatGPT and Claude is identified separately. No arXiv replacement or external correspondence is represented as completed.

\bibliographystyle{unsrtnat}
\bibliography{references}
\end{document}
